Consider a satellite observing x-ray sources, while orbiting the Earth in the
equatorial plane with orbital radius $r$, and orbital time period $P$. Let us
assume that this satellite is pointed to one fixed direction in space for a
given length of time. Take the radius of the earth as $R$.

When the satellite moves `behind' the earth, naturally, the x-ray source is
`occulted' and the measured x-ray flux from the source drops to zero. However,
due to Earth's atmosphere, this drop is gradual. If the line of sight of the
source passes through the atmosphere, the attenuation depends on the air-mass
(i.e.\ length of air column) along the line of sight.

\begin{description}
\item[(a)] Let us assume that pointing towards a fixed source at $0^\circ$
  declination. We consider that the source is occulted when 50\% of the light
  coming from the source gets attenuated due to the atmosphere. Let us say
  that this happens when the minimum height of the line of sight from the
  surface of the Earth is $h$.

  If $\theta_0$ is the angle between the direction to the source and the
  direction to the Earth, as measured from the spacecraft, find an expression
  for $\theta_0$. \hfill[1 point]
\item[(b)] The time duration $\Delta t$ between the source getting attenuated
  from 90\% of pre-occultation flux to 10\% is defined as the `occultation
  time' for the source. Assume the flux attenuates to 90\% when the minimum
  height of the line of sight is $(h + 0.5\Delta h)$ and similarly the flux
  attenuates to 10\% at $(h - 0.5\Delta h)$, where $\Delta h \ll R$.\\
  Find the expression for $\Delta t$ in terms of $r$, $P$, $\Delta h$ and
  $\theta_0$. \hfill[4 points]
\item[(c)] If the satellite was pointing towards a source at declination
  $\beta$ instead ($\beta$ not too large), what will be the expression for
  $\Delta t$? \hfill[15 points]
\end{description}

\textbf{Note:} If the satellite was not in the equatorial plane, then the
problem could have been simply rephrased by assuming the satellite's orbital
plane to be the equatorial plane. In that case, $\beta$ becomes `relative
declination'.
